% SHUORI visual tour, 10 October 2026.
% Original authored document/layout: CC BY-NC-SA-4.0.
% Screenshots contain software and creative assets under their respective terms.
% Build from this directory with XeLaTeX; see README.md.
\documentclass[10pt,a4paper]{article}
\usepackage[margin=16mm,top=19mm,bottom=17mm,headheight=20pt,headsep=5mm,footskip=8mm]{geometry}
\usepackage{fontspec}
\IfFontExistsTF{TeX Gyre Heros}{\setmainfont{TeX Gyre Heros}}{\setmainfont{Latin Modern Sans}}
\IfFileExists{fonts/NotoSansJP-Regular.ttf}
  {\newfontfamily\japanesefont[Path=fonts/]{NotoSansJP-Regular.ttf}}
  {\newfontfamily\japanesefont[Path=../../server/fonts/]{NotoSansJP-Regular.ttf}}
\usepackage{xcolor,graphicx,float,caption,fancyhdr,hyperref}
\definecolor{aqua}{HTML}{16799A}
\definecolor{navy}{HTML}{183B50}
\definecolor{muted}{HTML}{526878}
\definecolor{frame}{HTML}{D5E6ED}
\definecolor{pale}{HTML}{F0F8FB}
\hypersetup{
  pdftitle={SHUORI | A visual tour of volunteer operations},
  pdfauthor={HUANG WANHONG},
  pdfsubject={Twenty views of an independent hospital volunteer coordination application},
  pdfkeywords={SHUORI, volunteer coordination, scheduling, 3D, visual tour},
  colorlinks=true,urlcolor=aqua,linkcolor=aqua,pdfdisplaydoctitle=true,
  pdfpagelayout=SinglePage
}
\urlstyle{same}
\setlength{\parindent}{0pt}
\setlength{\parskip}{0pt}
\setlength{\intextsep}{0pt}
\setlength{\fboxsep}{0pt}
\setlength{\fboxrule}{0.4pt}
\emergencystretch=2em
\color{navy}
\DeclareCaptionFont{tour}{\fontsize{9}{11}\selectfont\color{muted}}
\DeclareCaptionLabelFormat{tour}{\textcolor{aqua}{\textbf{FIG.\ #2}}}
\captionsetup{font=tour,labelformat=tour,labelsep=quad,skip=3mm,
  width=176mm,justification=raggedright,singlelinecheck=false}
\pagestyle{fancy}
\fancyhf{}
\fancyhead[L]{\small\textcolor{aqua}{{\japanesefont 守織}\hspace{2mm}\textbf{SHUORI}}}
\fancyhead[R]{\fontsize{8}{10}\selectfont\textcolor{muted}{VOLUNTEER OPERATIONS / VISUAL TOUR}}
\fancyfoot[L]{\fontsize{7.5}{9}\selectfont\textcolor{muted}{Illustrative demo\enspace\textperiodcentered\enspace Fictional records}}
\fancyfoot[R]{\fontsize{8}{10}\selectfont\textcolor{aqua}{\textbf{\thepage}\ /\ 10}}
\renewcommand{\headrulewidth}{0.4pt}
\renewcommand{\footrulewidth}{0pt}
\renewcommand{\headrule}{\hbox to\headwidth{\color{frame}\leaders\hrule height\headrulewidth\hfill}}

% Every frame has the same 176 x 90 mm internal dimensions. Images are fitted,
% never stretched or cropped by this document; any UI crop is in the source PNG.
\newcommand{\tourimage}[1]{%
  \centering
  \fcolorbox{frame}{white}{%
    \parbox[c][90mm][c]{176mm}{\centering
      \includegraphics[width=176mm,height=90mm,keepaspectratio]{images/#1}}}%
}
\newcommand{\tourheading}[3]{%
  {\fontsize{8}{10}\selectfont\color{aqua}\textbf{#1}}\par
  \vspace{1.8mm}
  {\fontsize{22}{25}\selectfont\bfseries #2}\par
  \vspace{2.5mm}
  {\fontsize{9.5}{12}\selectfont\color{muted}#3}\par
  \vspace{5mm}
}
\newcommand{\figuregap}{\vspace{6mm}}

\begin{document}

\tourheading{01 / A CONNECTED WORKSPACE}{People, plans and place}{%
  An independent personal project for hospital volunteer coordination.\newline
  A visual tour of SHUORI 1.4.0, captured 10 October 2026 (JST). No organizational affiliation.}

\begin{figure}[H]
  \tourimage{01-overview.png}
  \caption{\textbf{Operations at a glance.} A selected-day overview brings staffing, service activity, upcoming events and project progress into one coordination workspace.}
  \label{fig:overview}
\end{figure}
\figuregap
\begin{figure}[H]
  \tourimage{02-building.png}
  \caption{\textbf{Explore the demonstration building.} An exploded 3D view connects eight modeled levels with service locations; each floor can be opened for detailed inspection.}
  \label{fig:building}
\end{figure}
\newpage

\tourheading{02 / VOLUNTEER RELATIONSHIPS}{A profile with a history}{%
  Keep current coordination information alongside dated follow-ups, qualifications and milestones.\newline
  Volunteer histories remain distinct from recorded service hours.}

\begin{figure}[H]
  \tourimage{03-volunteer-profile.png}
  \caption{\textbf{Rich volunteer profiles.} Contacts, programme status, skills, availability and readiness information support assignment decisions and individual follow-up. All displayed people are fictional.}
  \label{fig:volunteer-profile}
\end{figure}
\figuregap
\begin{figure}[H]
  \tourimage{04-volunteer-records.png}
  \caption{\textbf{A dated record of support.} Journal entries preserve training, conversations and recognition, with follow-up dates and optional links to scheduled events.}
  \label{fig:volunteer-records}
\end{figure}
\newpage

\tourheading{03 / FLEXIBLE INFORMATION}{Records that fit the work}{%
  Extend the information model with typed fields, then record actual service explicitly.\newline
  Saved definitions and values also accompany administrative Excel exports.}

\begin{figure}[H]
  \tourimage{05-custom-fields.png}
  \caption{\textbf{Configurable record fields.} An unsaved date-field definition illustrates how coordinators extend records. Text, number, choice and yes/no fields are also supported.}
  \label{fig:custom-fields}
\end{figure}
\figuregap
\begin{figure}[H]
  \tourimage{06-activity-records.png}
  \caption{\textbf{Account for completed service.} Activity records connect volunteers and shifts with actual hours and service counts, providing an explicit basis for monthly reporting.}
  \label{fig:activity-records}
\end{figure}
\newpage

\tourheading{04 / EVENTS AND MEETINGS}{An event becomes a workspace}{%
  Coordinate activities, training and meetings with configurable event types and linked modules.\newline
  Participants, preparation and operational references stay with the event.}

\begin{figure}[H]
  \tourimage{07-events.png}
  \caption{\textbf{A shared programme of events.} Searchable events combine date ranges, status, type and ownership; selected records can be exported as an agenda.}
  \label{fig:events}
\end{figure}
\figuregap
\begin{figure}[H]
  \tourimage{08-meeting.png}
  \caption{\textbf{Meeting information in context.} The configured platform and agenda stay with the event. This fictional example has no joining URL; meetings use external services.}
  \label{fig:meeting}
\end{figure}
\newpage

\tourheading{05 / PREPARATION AND HANDOVER}{Keep the details together}{%
  Place reports, shared-document links and preparation tasks where coordinators need them.\newline
  Local uploads use authenticated access; external services retain their own permissions.}

\begin{figure}[H]
  \tourimage{09-event-files.png}
  \caption{\textbf{Attach an event document.} The file/link editor accepts local uploads or an existing HTTPS address; this capture shows an unsaved example.}
  \label{fig:event-files}
\end{figure}
\figuregap
\begin{figure}[H]
  \tourimage{10-event-checklist.png}
  \caption{\textbf{Preparation made visible.} Completed and outstanding checklist items retain owners and due dates, making the next steps available within the event workspace.}
  \label{fig:event-checklist}
\end{figure}
\newpage

\tourheading{06 / DAILY AND WEEKLY SCHEDULING}{Make coverage readable}{%
  Examine assignments over time, then widen the view to coordinate the week.\newline
  Staffing suggestions and eligibility explanations support a coordinator's review.}

\begin{figure}[H]
  \tourimage{11-schedule.png}
  \caption{\textbf{The daily schedule.} Time-based assignments make service coverage and unfilled work visible, with scheduled events available alongside volunteer shifts.}
  \label{fig:schedule}
\end{figure}
\figuregap
\begin{figure}[H]
  \tourimage{12-weekly-schedule.png}
  \caption{\textbf{A wider planning horizon.} The weekly board groups activities by date, helping coordinators compare service commitments before assigning or adjusting work.}
  \label{fig:weekly-schedule}
\end{figure}
\newpage

\tourheading{07 / DELIVERY AND REFLECTION}{From dependencies to outcomes}{%
  Connect project preparation with operational delivery and recorded service results.\newline
  Planning measures and actual activity remain distinguishable throughout the workspace.}

\begin{figure}[H]
  \tourimage{13-project-plan.png}
  \caption{\textbf{A project plan with dependencies.} Task timing, progress and calculated critical-path information expose the activities that constrain the plan's completion.}
  \label{fig:project-plan}
\end{figure}
\figuregap
\begin{figure}[H]
  \tourimage{14-reports.png}
  \caption{\textbf{Review the month.} Service hours, participation and category trends summarize saved activity records and support monthly reports and recognition planning.}
  \label{fig:reports}
\end{figure}
\newpage

\tourheading{08 / DETAILED SPATIAL PLANNING}{A model you can inspect}{%
  Explore independently modeled furnishings, equipment, signs and volunteer figures.\newline
  Illustrative geometry provides a setting for discussion and layout studies.}

\begin{figure}[H]
  \tourimage{15-library-detail.png}
  \caption{\textbf{Inside a modeled service space.} Shelves with individual book spines, desks, screens and trolleys make the library tangible; selection reveals object information and transforms.}
  \label{fig:library-detail}
\end{figure}
\figuregap
\begin{figure}[H]
  \tourimage{16-asset-catalog.png}
  \caption{\textbf{An original asset library.} Fifteen procedural furnishing and equipment types can populate a floor study; each placed asset remains independently selectable.}
  \label{fig:asset-catalog}
\end{figure}
\newpage

\tourheading{09 / INTERACTION AND REHEARSAL}{Arrange, inspect, rehearse}{%
  Save layout studies and assignment-linked routes as named spatial scenarios.\newline
  Rehearsal animation illustrates a plan; it does not represent live tracking or validated navigation.}

\begin{figure}[H]
  \tourimage{17-object-inspector.png}
  \caption{\textbf{Inspect and position individual objects.} Selection connects a model object to its information and transform controls; editable furnishings can be moved or rotated.}
  \label{fig:object-inspector}
\end{figure}
\figuregap
\begin{figure}[H]
  \tourimage{18-route-rehearsal.png}
  \caption{\textbf{Draft a volunteer rehearsal route.} Four unsaved waypoints connect a scheduled person to a schematic floor path. Routes do not detect obstacles or verify access.}
  \label{fig:route-rehearsal}
\end{figure}
\newpage

\tourheading{10 / DOCUMENTS FOR COORDINATION}{Turn a plan into a handout}{%
  Preview and export volunteer timelines, station coverage, weekly rosters, agendas and event briefs.\newline
  Choose PDF or editable Excel, with A4 and A3 layouts.}

\begin{figure}[H]
  \tourimage{19-export-studio.png}
  \caption{\textbf{An export studio for operational views.} Date, time and status controls define the saved-data scope; a preview supports review before downloading.}
  \label{fig:export-studio}
\end{figure}
\figuregap
\begin{figure}[H]
  \tourimage{20-event-brief.png}
  \caption{\textbf{Prepare a readable event brief.} Export event details, preparation, participants and linked operational information, with coordination notes and custom fields available by choice.}
  \label{fig:event-brief}
\end{figure}
\vspace{4mm}
{\fontsize{6.9}{8.5}\selectfont\color{muted}
  \textbf{Provenance.} Demo floor geometry was independently reconstructed from the
  \href{https://www.h.u-tokyo.ac.jp/english/international-patients/floor-guide/index.html}{University of Tokyo Hospital's public floor-guide maps},
  with invented furnishings and routes. It is illustrative, not official, surveyed or suitable for navigation.
  No hospital or organizational affiliation is implied.\par
  \vspace{1mm}
  \textcopyright\ 2026 HUANG WANHONG and contributors.
  \href{https://github.com/wwwwanhonghuang/Hospital-Volunteer-Activity-Management}{Project and source notices}:
  software AGPL-3.0-only; original creative assets and authored document CC BY-NC-SA 4.0.
  Third-party components retain their terms; Noto Sans JP uses SIL OFL 1.1.
}

\end{document}
